\begin{table}[H]
\centering
\caption*{\textbf{Quadro 7 - Elementos removidos da especificação mínima}}
\scriptsize
\renewcommand{\arraystretch}{1.12}
\begin{tabularx}{\linewidth}{|p{0.34\linewidth}|Y|}
\hline
\textbf{Elemento removido} & \textbf{Razão} \\
\hline
Extração obrigatória de hints antes de P1 &
Nos diagnósticos, a compressão perdeu sinal procedural; a 0.3 filtra candidatas e
preserva os bodies canônicos das skills mantidas. \\
\hline
Compilador por LLM com muitos campos inferidos &
\texttt{name}, \texttt{description} e body já sustentam indexação e reranking. \\
\hline
Tool router independente &
\(T_{\mathrm{base}}\) e a união de \texttt{allowed-tools} resolvem o menu mínimo. \\
\hline
Camada de capabilities e providers &
Nomes diretos de tools bastam enquanto não houver implementações equivalentes. \\
\hline
Resolvedor de \texttt{scripts/}, \texttt{references/} e \texttt{assets/} &
O conceito é expresso carregando apenas \texttt{SKILL.md}. \\
\hline
Otimizador formal de bundles &
Seleção ordenada, limite e critérios qualitativos são suficientes. \\
\hline
Linguagem formal de precedência &
A ordem estável do reranking é a convenção da MOSAIC 0.3. \\
\hline
Verificador separado &
\texttt{doneWhen} pode ser avaliado pelo executor na especificação mínima. \\
\hline
Protocolo completo de erros &
Status de nó, limites e pedido semântico de revisão descrevem o fluxo central. \\
\hline
IDs do provedor em saídas do modelo &
Correlação pertence ao runtime; o modelo não deve copiar referências opacas. \\
\hline
Segurança, autorização e prompt injection &
A especificação assume catálogo, tools e runtime confiáveis em modo YOLO. \\
\hline
Scheduler distribuído e paralelismo formal &
Um DAG expressa a ordem necessária sem infraestrutura adicional. \\
\hline
\end{tabularx}
\par\smallskip\raggedright\small Fonte: elaboração própria (2026).
\end{table}
